% Quick Start, Chapter 4 -- Formulas and @Functions.  Figure screens/qs4.k123.
% Values replayed through the shipped application (examples/ops4.i23).
% Sources (Law 1): os/i123/core/formula.c (operators and their precedence),
% os/i123/core/eval.c (@SUM, @AVG, @MAX, @MIN, @COUNT, @IF, @ROUND),
% os/i123/app/face_a.c (/Copy).
% GATES: test-i123-formula test-pubs-i123-examples test-pubs-screens test-pubs-i123
\Gates{test-i123-formula, test-pubs-i123-examples, test-pubs-screens,
  test-pubs-i123}
\chapter{Formulas and @Functions}

This chapter prices three supplies and summarizes the costs with
functions.  Start with an empty worksheet.

\section{What This Chapter Shows You}
\begin{itemize}
\item Writing a formula with operators
\item Copying a formula down a column
\item Summarizing a range with \fn{@SUM}, \fn{@AVG} and \fn{@MAX}
\end{itemize}

\section{Formulas}
\ix{formulas}\ix{operators}
A formula that begins with a cell address needs a plus sign in front, so
that Initech 123 reads it as a value and not as a label.  The operators are
\texttt{+}, \texttt{-}, \texttt{*}, \texttt{/} and \texttt{\^{}} (power).
Exponents come first, then multiplication and division, then addition and
subtraction; parentheses change the order.
\LotExample{examples/ops4}

\section{Entering the Supplies}
\ix{"@AVG}\ix{"@MAX}
\begin{Keystrokes}
\Do{Type}{\type{Pens}}{in A1; in B1 and C1 type \type{12} and \type{.5},
  pressing \kname{RIGHT} between the entries}
\Do{Type}{\type{+b1*c1}}{in D1, the quantity times the price, and press
  \kname{ENTER}}
\Do{Type}{\type{Pads}}{in A2, \type{8} in B2 and \type{1.25} in C2}
\Do{Type}{\type{Tape}}{in A3, \type{6} in B3 and \type{.75} in C3}
\Do{Press}{\kname{SLASH}}{type \type{c}, type \type{d1} and press
  \kname{ENTER}, type \type{d2} and press \kname{ENTER}}
\Do{Press}{\kname{SLASH}}{type \type{c}, type \type{d1} and press
  \kname{ENTER}, type \type{d3} and press \kname{ENTER}}
\end{Keystrokes}
The copies refer to their own rows: D2 holds \texttt{+B2*C2}.

\section{Summarizing with Functions}
\begin{Keystrokes}
\Do{Press}{\fkey{Goto}{F5}}{type \type{d5} and press \kname{ENTER}}
\Do{Type}{\type{@sum(d1..d3)}}{and press \kname{DOWN}}
\Do{Type}{\type{@avg(d1..d3)}}{and press \kname{DOWN}}
\Do{Type}{\type{@max(d1..d3)}}{and press \kname{ENTER}}
\end{Keystrokes}
\FlairShot[fig:qs4]{qs4}{Costs, with their total, average and largest}

\section{Task Summary}
\begin{LotTable}{p{58mm}p{70mm}}
\LotHead{\textbf{To} & \textbf{You used}}
Multiply two cells & \texttt{+b1*c1}\\
Copy a formula to the next row & /Copy, the cell to copy and the cell below\\
Add a range & \fn{@SUM}\\
Find the average & \fn{@AVG}\\
Find the largest value & \fn{@MAX}\LotEnd
\end{LotTable}

\section{Other Tasks to Try}
\begin{LotTable}{p{58mm}p{70mm}}
\LotHead{\textbf{To} & \textbf{Use}}
Find the smallest cost & \texttt{@MIN(D1..D3)}\\
Count the entries & \texttt{@COUNT(D1..D3)}\\
Round the average to cents & \texttt{@ROUND(D6,2)}\\
Mark the expensive items & \texttt{@IF(D1>8,"High","Low")}\LotEnd
\end{LotTable}
